\documentclass[11pt,a4paper]{article}

\usepackage[utf8]{inputenc}
\usepackage[T1]{fontenc}
\usepackage[spanish,es-nodecimaldot]{babel}
\usepackage[margin=2.4cm]{geometry}
\usepackage{booktabs}
\usepackage{longtable}
\usepackage{array}
\usepackage{xcolor}
\usepackage{colortbl}
\usepackage[hidelinks]{hyperref}
\usepackage{enumitem}
\usepackage{fancyhdr}
\usepackage{titlesec}

\definecolor{tinta}{HTML}{1B2A41}
\definecolor{acento}{HTML}{C75B12}
\definecolor{suave}{HTML}{F2F4F7}

\titleformat{\section}{\Large\bfseries\color{tinta}}{\thesection.}{0.6em}{}
\titleformat{\subsection}{\large\bfseries\color{tinta}}{\thesubsection.}{0.6em}{}

\pagestyle{fancy}
\fancyhf{}
\lhead{\small\color{tinta}CUMBRE — Análisis de mercado y estructura de precios}
\rfoot{\small\thepage}
\renewcommand{\headrulewidth}{0.4pt}

\newcommand{\aviso}[1]{%
  \begin{center}
  \fcolorbox{acento}{suave}{%
    \begin{minipage}{0.93\textwidth}
    \vspace{4pt}#1\vspace{4pt}
    \end{minipage}}
  \end{center}}

\begin{document}

\begin{center}
{\Huge\bfseries\color{tinta} CUMBRE}\\[4pt]
{\Large\color{tinta} Análisis de mercado y estructura de precios}\\[10pt]
{\small Plataforma de aprendizaje adaptativo — Vastoria Solutions S.A.C.}\\
{\small 29 de julio de 2026}
\end{center}

\vspace{10pt}

\section{Qué es y qué no es este documento}

\aviso{\textbf{Ninguna plataforma del mercado peruano publica sus precios.}
Todas —sin excepción— exigen contacto comercial para cotizar. Se verificó
también el mercado colombiano, con el mismo resultado. Por tanto, este
documento \textbf{no contiene una tabla de precios de la competencia}: sería
inventada.}

Lo que sí contiene, y está verificado:

\begin{itemize}[leftmargin=1.3em]
  \item \textbf{Cómo empaqueta la competencia.} La estructura por módulos sí
        es pública y se documenta más abajo.
  \item \textbf{Qué modelo de cobro usan.} Declarado por los propios
        proveedores en sus sitios.
  \item \textbf{Cuánto cuesta operar CUMBRE.} Calculado sobre la
        configuración real desplegada, no sobre supuestos.
  \item \textbf{Qué estructura de precios se sostiene} con esos costos.
\end{itemize}

Las cifras marcadas como \emph{estimación} son órdenes de magnitud sobre
precios de lista públicos. Deben confirmarse contra la facturación real antes
de usarse en una negociación.

\section{El mercado peruano}

\subsection{Quién está y cómo se posiciona}

\begin{longtable}{@{}p{3.1cm}p{5.2cm}p{6.3cm}@{}}
\toprule
\textbf{Plataforma} & \textbf{Posicionamiento} & \textbf{Señal relevante} \\
\midrule
\endhead
Smiledu & Gestión educativa por paquetes & Vende en bloques de \textbf{8 y 12
módulos}; tesorería y asistencia como adicionales \\
\addlinespace
Cubicol & Volumen y precio accesible & Más de \textbf{800 instituciones} en
Perú; apps móviles nativas \\
\addlinespace
SophIA & Cumplimiento y confianza & \textbf{Autorizada por SUNAT} para boletas
electrónicas; \textbf{ISO 27001}; más de 10 años \\
\addlinespace
JSedu & Flexibilidad de plan & Cobra por \textbf{número de alumnos y módulos}
contratados \\
\addlinespace
SIGEDU & Cobertura por tamaño & Desde centros pequeños hasta grandes
instituciones \\
\addlinespace
DataCole & ERP académica & Enfoque administrativo integral \\
\addlinespace
SIADED & Retención & Módulos nuevos y actualizaciones \textbf{sin costo
adicional} \\
\bottomrule
\end{longtable}

\subsection{Modelos de cobro observados}

\begin{longtable}{@{}p{4.2cm}p{10.4cm}@{}}
\toprule
\textbf{Eje} & \textbf{Cómo se aplica} \\
\midrule
\endhead
Por alumno & Cuota mensual proporcional a la matrícula del centro. Es el eje
dominante. \\
\addlinespace
Por módulo & Un núcleo obligatorio más complementos contratables por separado.
\\
\addlinespace
Por escalón & Planes cerrados (básico / operación / crecimiento) en lugar de
precio unitario. \\
\addlinespace
Implementación & Costo aparte de puesta en marcha y capacitación, señalado en
el sector como partida frecuentemente olvidada al comparar. \\
\bottomrule
\end{longtable}

\aviso{\textbf{Dato estratégico:} que nadie publique precios significa que en
este mercado \textbf{se vende con reunión, no con botón de compra}. El precio
se construye por cliente. Eso favorece a quien puede justificar su costo con
números —lo que hace útil la sección siguiente.}

\section{Costo real de operar CUMBRE}

Calculado sobre la infraestructura efectivamente desplegada al 29 de julio de
2026.

\subsection{Configuración actual}

\begin{longtable}{@{}p{4.2cm}p{10.4cm}@{}}
\toprule
\textbf{Componente} & \textbf{Configuración real} \\
\midrule
\endhead
Cloud Run & 3 servicios (auth, content, learning). 512\,MiB, mínimo
\textbf{1 instancia siempre viva}, máximo 10. Región
\texttt{southamerica-west1}. \\
\addlinespace
Cloud SQL & PostgreSQL 16, edición \emph{enterprise}, nivel
\texttt{db-f1-micro}, disco 10\,GB con crecimiento automático. \\
\addlinespace
Vercel & 4 proyectos web. Hoy en plan Hobby. \\
\addlinespace
Modelo de lenguaje & \texttt{gpt-4o-mini} vía OpenAI, por consumo. \\
\addlinespace
Dominios & Bajo \texttt{teamvastoria.com}, ya adquirido. \\
\bottomrule
\end{longtable}

\subsection{Estructura de costo mensual}

\begin{longtable}{@{}p{3.6cm}p{4.6cm}p{6.4cm}@{}}
\toprule
\textbf{Partida} & \textbf{Cómo se cobra} & \textbf{Comentario} \\
\midrule
\endhead
Cloud Run & Por CPU/memoria asignada y peticiones & La instancia mínima
siempre viva es un \textbf{costo fijo}, no proporcional al uso. Se eligió a
propósito: sin ella, el primer alumno de la mañana espera varios segundos. \\
\addlinespace
Cloud SQL & Fijo por nivel + disco & \texttt{db-f1-micro} es el más pequeño.
Soporta varias academias antes de necesitar cambio. \\
\addlinespace
Vercel Pro & Fijo por plan & \textbf{Obligatorio al firmar}: el plan Hobby es
para uso no comercial. \\
\addlinespace
OpenAI & Por token consumido & \textbf{Única partida que escala con el uso}.
Crece con alumnos activos, no con alumnos matriculados. \\
\bottomrule
\end{longtable}

\aviso{\textbf{Antes de fijar precio, obtén el número real.} La consola de
facturación de Google Cloud ya tiene semanas de datos del proyecto
\texttt{vastoria-campus}. Ese número es el punto de partida honesto; cualquier
estimación de lista es inferior a él en precisión.}

\subsection{Lo que la estructura de costo implica para el precio}

\begin{enumerate}[leftmargin=1.4em]
  \item \textbf{Tu costo es mayoritariamente fijo.} Infraestructura siempre
        encendida más planes mensuales. Eso significa que el primer cliente es
        caro y el segundo casi gratis: \emph{el margen mejora con cada
        academia, no con cada alumno}.
  \item \textbf{La única partida variable es la IA.} Y depende de alumnos
        \emph{activos}, no matriculados. Un colegio de 500 alumnos donde 50
        usan el tutor cuesta como uno de 50.
  \item \textbf{Conclusión de negocio:} cobrar por alumno matriculado —el
        modelo dominante del mercado— te da ingresos que crecen más rápido que
        tus costos. Es el modelo correcto para ti, y además es el que el
        cliente ya entiende porque es el de todos.
\end{enumerate}

\section{Estructura de precios propuesta}

Se propone replicar el modelo del mercado —núcleo más módulos— porque es el
que el comprador ya sabe evaluar.

\begin{longtable}{@{}p{3.8cm}p{5.4cm}p{5.4cm}@{}}
\toprule
\textbf{Bloque} & \textbf{Qué incluye} & \textbf{Estado hoy} \\
\midrule
\endhead
\rowcolor{suave}
Campus base & Acceso por rol, aulas, temas, lecciones, progreso, panel de
dirección & \textbf{Funcionando} \\
\addlinespace
Tutor con IA & Tutor conversacional anclado a la lección, con contexto del
progreso del alumno & \textbf{Funcionando}; requiere clave de OpenAI activa \\
\addlinespace
Ruta adaptativa & Secuenciación por dominio y grafo de conocimiento &
Funcionando; enriquecimiento del grafo pendiente de corregir \\
\addlinespace
Autoría y generación & Constructor de módulos y generación asistida &
Parcial \\
\addlinespace
YariNET & Retos & \textbf{En construcción}; oculto para clientes \\
\addlinespace
Marca propia & Dominio, logo y colores de la institución & \textbf{Funcionando}
\\
\bottomrule
\end{longtable}

\aviso{\textbf{El diferenciador real:} las siete plataformas del cuadro
anterior son sistemas de \emph{gestión} —matrícula, notas, pensiones,
asistencia—. CUMBRE es una plataforma de \emph{aprendizaje}. No compiten en lo
mismo.\\[4pt]
Eso es ventaja y riesgo a la vez: no hay con quién compararte, pero el colegio
ya le paga a una de ellas y hay que explicarle por qué necesita otra cosa. La
venta no es «somos mejores que Cubicol», es «esto hace algo que Cubicol no
hace».}

\section{Barreras de entrada que hoy no tienes}

\begin{longtable}{@{}p{4.2cm}p{10.4cm}@{}}
\toprule
\textbf{Barrera} & \textbf{Situación} \\
\midrule
\endhead
Autorización SUNAT & SophIA la tiene para emisión de boletas electrónicas. Si
algún día se cobran pensiones desde la plataforma, es requisito, no adorno. \\
\addlinespace
ISO 27001 & Certificación de seguridad de la información. Pesa en licitaciones
y en colegios grandes. \\
\addlinespace
Base instalada & Cubicol supera las 800 instituciones. La referencia de pares
vende sola en este sector. \\
\addlinespace
App móvil nativa & Varios competidores la ofrecen. CUMBRE es web
adaptable. \\
\bottomrule
\end{longtable}

\section{Recomendaciones}

\begin{enumerate}[leftmargin=1.4em]
  \item \textbf{Obtener el costo real} desde la consola de facturación antes
        de cualquier propuesta económica.
  \item \textbf{Cobrar por alumno matriculado}, con núcleo más módulos:
        alineado al mercado y favorable a tu estructura de costos.
  \item \textbf{Cotizar la implementación aparte.} El sector lo hace y el
        cliente lo espera; incluirlo silenciosamente erosiona el margen.
  \item \textbf{No competir contra los sistemas de gestión.} Posicionar como
        complemento pedagógico, no como reemplazo administrativo.
  \item \textbf{Pasar a Vercel Pro al firmar.} El plan actual no permite uso
        comercial.
\end{enumerate}

\section*{Fuentes}

\small
\begin{itemize}[leftmargin=1.3em]
  \item Smiledu — \url{https://smiledu.com/}
  \item Cubicol — \url{https://www.cubicol.pe/}
  \item JSedu — \url{https://www.jsedu.company/sistema-gestion-escolar-peru}
  \item SIGEDU — \url{https://www.sigedu.pe/}
  \item DataCole — \url{https://www.datacole.com/}
  \item SIADED — \url{https://siaded.com.pe/}
  \item Tecnomedia, comparativa de software de gestión escolar en Perú —
        \url{https://www.tecnomedia.pe/comparaciones/mejores-software-gestion-escolar-peru/}
  \item Edtools, software para colegios en Colombia —
        \url{https://edtools.co/blog/software-para-colegios-colombia-precios-modulos.html}
\end{itemize}

\vspace{6pt}
\noindent\textit{Configuración de infraestructura tomada de
\texttt{deploy/gcp/desplegar.sh} y \texttt{deploy/gcp/preparar.sh} del
repositorio, no de estimaciones.}

\end{document}
